% BEGIN VOL08-EXPANSION VIII10
\section{Supplementary worked examples}
\label{sec:viii10-pedagogy}

\begin{example}[Exponential covering of the circle]\label{ex:viii10-ped-01}
The map $p:\mathbb R\to S^1$, $p(t)=e^{2\pi it}$, evenly covers every sufficiently short arc.  Each fiber is $t+\mathbb Z$, making $\mathbb R$ the universal cover of $S^1$.
\end{example}

\begin{example}[Power covering]\label{ex:viii10-ped-02}
The map $p_n:S^1\to S^1$, $p_n(z)=z^n$, is an $n$-sheeted covering.  A small arc avoiding wraparound has $n$ disjoint inverse arcs, each mapped homeomorphically onto it.
\end{example}

\begin{example}[Product covering]\label{ex:viii10-ped-03}
If $p:E\to B$ is a covering and $X$ is any space, then $p\times\operatorname{id}_X:E\times X\to B\times X$ is a covering; evenly covered neighborhoods are products of evenly covered opens with opens in $X$.
\end{example}

\section{Graded supplementary exercises}
The following exercises progress from direct computation and construction to proof, hypothesis testing, application, and synthesis.

\subsection*{Standard computations and constructions}

\begin{exercise}[Evenly covered]\label{exr:viii10-ped-01}
State what it means for an open set $U\subset B$ to be evenly covered by $p:E\to B$.
\end{exercise}
\begin{hint}
Describe $p^{-1}(U)$ as disjoint sheets.
\end{hint}
\begin{solution}
$p^{-1}(U)$ is a disjoint union of open sets $V_\alpha$, each mapped homeomorphically onto $U$ by $p$.
\end{solution}

\begin{exercise}[Fiber of exponential cover]\label{exr:viii10-ped-02}
Compute $p^{-1}(1)$ for $p(t)=e^{2\pi it}$.
\end{exercise}
\begin{hint}
Solve $e^{2\pi it}=1$.
\end{hint}
\begin{solution}
The fiber is $\mathbb Z$.
\end{solution}

\begin{exercise}[Fiber of power map]\label{exr:viii10-ped-03}
How many points lie over $1$ under $z\mapsto z^n$?
\end{exercise}
\begin{hint}
Solve $z^n=1$.
\end{hint}
\begin{solution}
Exactly $n$, the $n$th roots of unity.
\end{solution}

\begin{exercise}[Local homeomorphism]\label{exr:viii10-ped-04}
Why is every covering map a local homeomorphism?
\end{exercise}
\begin{hint}
Use one sheet of an evenly covered neighborhood.
\end{hint}
\begin{solution}
For $e\in E$, choose an evenly covered $U$ around $p(e)$; the sheet containing $e$ maps homeomorphically to $U$.
\end{solution}

\begin{exercise}[Product sheets]\label{exr:viii10-ped-05}
If $U$ is evenly covered by $p$ and $W\subset X$ is open, describe the sheets over $U\times W$ for $p\times\operatorname{id}_X$.
\end{exercise}
\begin{hint}
Take products with each sheet.
\end{hint}
\begin{solution}
They are $V_\alpha\times W$, each homeomorphic to $U\times W$.
\end{solution}

\subsection*{Proofs}

\begin{exercise}[Discrete fibers]\label{exr:viii10-ped-06}
Prove fibers of a covering map are discrete subspaces of $E$.
\end{exercise}
\begin{hint}
Use an evenly covered neighborhood around the image point.
\end{hint}
\begin{solution}
Each sheet contains exactly one point of a given fiber, and the sheet is an open neighborhood separating that point from the other fiber points.
\end{solution}

\begin{exercise}[Constant sheet number]\label{exr:viii10-ped-07}
Prove a covering over a connected base has constant fiber cardinality.
\end{exercise}
\begin{hint}
The number of sheets is locally constant.
\end{hint}
\begin{solution}
An evenly covered neighborhood identifies fibers over all its points with the same set of sheets; hence the cardinality is locally constant and therefore constant on a connected base.
\end{solution}

\begin{exercise}[Composition]\label{exr:viii10-ped-08}
Under standard hypotheses, prove a composition of covering maps is a covering.
\end{exercise}
\begin{hint}
Refine evenly covered neighborhoods successively.
\end{hint}
\begin{solution}
Choose an evenly covered neighborhood for the lower cover and then neighborhoods whose sheets are evenly covered by the upper cover; a common refinement produces disjoint composite sheets.
\end{solution}

\begin{exercise}[Pullback preview]\label{exr:viii10-ped-09}
Explain why pulling back a covering along a map again yields a covering.
\end{exercise}
\begin{hint}
Pull back each local trivialization.
\end{hint}
\begin{solution}
Over an evenly covered $U$, the cover is a disjoint union of copies of $U$; its fiber product with $X\to B$ is a disjoint union of copies of the preimage of $U$, giving local covering charts.
\end{solution}

\subsection*{Counterexamples and hypothesis tests}

\begin{exercise}[Local homeomorphism is covering]\label{exr:viii10-ped-10}
Test: every local homeomorphism is a covering map.
\end{exercise}
\begin{hint}
Covering requires simultaneous evenly covered neighborhoods for all points of a fiber.
\end{hint}
\begin{solution}
False without additional hypotheses; local homeomorphism is a pointwise condition.
\end{solution}

\begin{exercise}[Fibers connected]\label{exr:viii10-ped-11}
Test: fibers of a nontrivial covering are connected.
\end{exercise}
\begin{hint}
Fibers are discrete.
\end{hint}
\begin{solution}
False; a multi-sheeted fiber is a discrete set with several points.
\end{solution}

\begin{exercise}[Sheet count can jump]\label{exr:viii10-ped-12}
Test: an $n$-sheeted covering over a connected base may have different fiber sizes at different points.
\end{exercise}
\begin{hint}
Use local constancy.
\end{hint}
\begin{solution}
False. The sheet number is constant on each connected component of the base.
\end{solution}

\subsection*{Applications and investigations}

\begin{exercise}[Unwrapping periodicity]\label{exr:viii10-ped-13}
Explain how the universal cover $\mathbb R\to S^1$ converts angular periodicity into translation symmetry.
\end{exercise}
\begin{hint}
Angles differing by integers project to the same point.
\end{hint}
\begin{solution}
Upstairs, the circle coordinate becomes a real coordinate; the lost periodicity reappears as the deck translations $t\mapsto t+k$.
\end{solution}

\begin{exercise}[Covering as local copies]\label{exr:viii10-ped-14}
Why are covering spaces useful for transferring local calculations?
\end{exercise}
\begin{hint}
Each sheet is locally identical to the base.
\end{hint}
\begin{solution}
Any local construction on an evenly covered neighborhood can be lifted independently to each sheet, while global topology is encoded by how the sheets connect and permute.
\end{solution}

\subsection*{Challenge problems}

\begin{exercise}[Universal-cover criterion preview]\label{exr:viii10-ped-15}
What extra global property distinguishes a universal cover from an arbitrary connected cover?
\end{exercise}
\begin{hint}
Its total space has trivial fundamental group.
\end{hint}
\begin{solution}
A universal cover is connected and simply connected, so every other connected cover is obtained as a quotient under the standard classification hypotheses.
\end{solution}

\begin{exercise}[Circle-cover synthesis]\label{exr:viii10-ped-16}
Relate the $n$-fold cover $S^1\to S^1$ to the universal cover $\mathbb R\to S^1$.
\end{exercise}
\begin{hint}
Quotient the universal cover by a subgroup of integer translations.
\end{hint}
\begin{solution}
The subgroup $n\mathbb Z$ acts on $\mathbb R$ by translations; $\mathbb R/n\mathbb Z\cong S^1$, and the induced map to $\mathbb R/\mathbb Z\cong S^1$ is the $n$-fold cover.
\end{solution}

% END VOL08-EXPANSION VIII10
